\documentclass[a4paper,12pt]{article}

\usepackage{alltt, fancyvrb, url}
\usepackage{graphicx}
\usepackage[utf8]{inputenc}
\usepackage{float}
\usepackage{hyperref}
\usepackage[italian]{babel}
\usepackage[table]{xcolor}
\usepackage{array}
\usepackage{multirow}
\usepackage{titlesec}
\usepackage{listings}
\usepackage[italian]{cleveref}
\usepackage{acronym}

\renewcommand{\thesection}{\arabic{section}}
\renewcommand{\labelenumii}{\arabic{enumi}.\arabic{enumii}}
\renewcommand{\labelenumiii}{\arabic{enumi}.\arabic{enumii}.\arabic{enumiii}}
\renewcommand{\labelenumiv}{\arabic{enumi}.\arabic{enumii}.\arabic{enumiii}.\arabic{enumiv}}

\sloppy
\setlength{\parindent}{0pt}

\title{\textbf{Requirements Breakdown Structure (RBS)}}
\author{}
\date{}

\begin{document}
\maketitle

\textbf{Goal:} entro un anno sviluppare un'applicazione software per gestire le colonnine di ricarica dell'azienda in
modo da ottenere un aumento del fatturato del 10\% entro l'anno successivo.
\begin{enumerate}
    \item Fornire agli utenti un'esperienza personalizzata in base al proprio profilo;
    \begin{enumerate}
        \item registrazione di un nuovo utente;
        \begin{enumerate}
            \item creazione di un database per gli account degli utenti;
            \item adozione di misure di sicurezza per proteggere i dati degli utenti;
        \end{enumerate}
        \item login di un utente registrato;
        \item registrazione e gestione di auto da parte degli utenti;
        \item memorizzazione delle impostazioni personalizzate;
        \item modifica dati dell'account.
    \end{enumerate}
    \item mostrare agli utenti le colonnine di ricarica presenti entro un raggio variabile rispetto alla posizione
    dell'utente e la loro disponibilità in tempo reale;
    \begin{enumerate}
        \item integrazione con il database contenente le posizioni delle colonnine di ricarica;
        \item visualizzazione delle colonnine di ricarica, differenziandole in base allo stato, su una mappa
        interattiva centrata sulla posizione dell'utente;
        \item aggiornamento del software delle colonnine di ricarica per inviare aggiornamenti in tempo reale sulla
        loro disponibilità;
    \end{enumerate}
    \item dare agli utenti la possibilità di avviare/interrompere una ricarica presso una colonnina di una propria auto
    registrata;
    \begin{enumerate}
        \item aggiunta della nozione di ricarica nei dati dei database;
        \item aggiornamento del software delle colonnine di ricarica per associare una ricarica ad un'auto registrata
        con l'applicazione;
        \item aggiunta funzionalità per avviare una ricarica associando un'auto registrata con una colonnina di ricarica;
    \end{enumerate}
    \item permettere agli utenti di pianificare un itinerario di viaggio per ottenere un percorso comprendente le soste
    di ricarica necessarie.
    \begin{enumerate}
        \item generazione di un percorso ottimizzato rispetto al tempo di percorrenza, che consideri le soste di ricarica
        necessarie in base alla batteria dell'auto utilizzata per il viaggio.
    \end{enumerate}
\end{enumerate}

\section*{Rischi, ostacoli ed assunzioni}
Relativamente alla numerazione dei requisiti, si possono individuare i seguenti rischi, ostacoli ed assunzioni:
\begin{itemize}
    \item 2.1 - \textbf{assunzione}: esiste già un database contenente le informazioni relative alle colonnine di ricarica;
    \item 2.3, 3.2 - \textbf{rischio}: introduzione di bug nel software delle colonnine di ricarica;
    \item 4.1 - \textbf{rischio}: incertezza nella realizzazione di un algoritmo per generare il percorso ottimizzato.
\end{itemize}

\end{document}